\clearpage\thispagestyle{empty}
\pdfbookmark[0]{Quellen zur Besprechung von Bodes Jahrbuch für 1810}{bode1810review-sources}
\begin{center}\textbf{Quellen zur Besprechung von Bodes Jahrbuch für 1810}\end{center}
\noindent Diese KI-Übertragung ins heutige Deutsch folgt Gauß’ Besprechung auf S.~504--507 der \href{https://gdz.sub.uni-goettingen.de/id/PPN236007467}{\textit{Werke}, Bd.~VI} (Göttingen, 1874). Sie erschien ohne Namensnennung unter der Überschrift \emph{Berlin} in den \textit{Göttingischen gelehrten Anzeigen}, Stück~155, 26. September 1808, S.~\href{https://api.digitale-sammlungen.de/iiif/image/v2/bsb10538674_00515/full/1800,/0/default.jpg}{1545}--\href{https://api.digitale-sammlungen.de/iiif/image/v2/bsb10538674_00522/full/1800,/0/default.jpg}{1552}. Das verlinkte Exemplar stammt aus der Bayerischen Staatsbibliothek.
\par\smallskip\noindent Gauß bespricht Johann Elert Bodes \emph{Astronomisches Jahrbuch für das Jahr} 1810, \emph{nebst einer Sammlung der neuesten in die astronomischen Wissenschaften einschlagenden Abhandlungen und Nachrichten}, 1807 in Berlin gedruckt. Das Jahr 1810 bezeichnet den Jahrgang.
\par\smallskip\noindent Der \href{https://api.digitale-sammlungen.de/iiif/image/v2/bsb10538674_00521/full/1800,/0/default.jpg}{Erstdruck} nennt bei der Sternbedeckung unter Nummer 13 den 20. März und gibt den Beobachtungen an der Berliner Sternwarte die Nummer 24. In der Werkausgabe steht im Datum ein c-förmiges Zeichen. Der Eintrag zur Sternwarte steht dort ohne Nummer. Die Übertragung folgt der Werkausgabe. KI-Übertragung und KI-Satz.
\clearpage
